%! Author = WelJo
%! Date = 8/18/2026

\newpage

\section{Fazit und Ausblick}
\label{sec:conclusion}

% - Beantwortung der Forschungsfrage

This study extended the Reliability of Accessibility (RoA) index to capture interdependencies between urban transportation and utility infrastructure during extreme floods.
Integrating a four-state finite state machine with road-network-constrained coupling advances evaluation beyond simple inundation models to full functional resilience.
Findings from the Trier case study demonstrate that:
\begin{enumerate}
    \item \textbf{Cascading failures} significantly amplify accessibility loss, frequently impacting areas far beyond the immediate flood zone.
    \item \textbf{Internal facility buffers and reboot delays} create critical temporal lags in service restoration.
    \item \textbf{A multi-tier evaluation approach} is essential to isolate specific network vulnerabilities and compound risks.
\end{enumerate}

\noindent
Future work should incorporate dynamic utility routing, probabilistic failure modeling, and real-time population movement data.
Ultimately, this framework equips planners and emergency managers to quantify service resilience and prioritize targeted infrastructure hardening against climate hazards.